% Generated draft: requires mathematical and source review.
\documentclass[11pt]{article}
\usepackage[margin=25mm]{geometry}
\usepackage{fontspec}
\setmainfont{Latin Modern Roman}
\usepackage{amsmath,amssymb,mathtools,unicode-math}
\setmathfont{Latin Modern Math}
\usepackage{enumitem,needspace,xurl,hyperref}
\setlength{\parindent}{0pt}
\setlength{\parskip}{5pt}
\setlength{\emergencystretch}{3em}
\newcommand{\N}{\mathbb{N}}
\newcommand{\Z}{\mathbb{Z}}
\newcommand{\Q}{\mathbb{Q}}
\newcommand{\R}{\mathbb{R}}
\newcommand{\C}{\mathbb{C}}
\newcommand{\F}{\mathbb{F}}
\newcommand{\A}{\mathbb{A}}
\newcommand{\PP}{\mathbb{P}}
\newcommand{\E}{\mathbb{E}}
\newcommand{\eps}{\varepsilon}
\newcommand{\dd}{\,\mathrm{d}}
\begin{document}
\section*{UnsolvedMath ID: 30005294}
\textbf{Source problem code:} catalogue entry
\subsection{Definitions and mathematical statement}
Let \(s\in(0,1)\) and let \(k\geq2\) be an integer. For a nonempty open set \(\Omega\subset\R^2\), define its inradius by
\[
r_\Omega:=\sup\{r>0:\text{there exists }x\in\Omega\text{ with }B_r(x)\subset\Omega\},
\qquad B_r(x):=\{y\in\R^2:|y-x|<r\}.
\]
The original source excerpt requires \(r_\Omega<\infty\); this restriction must be retained even though it is absent from the supplied clean statement. No boundedness or boundary regularity assumption is stated.

Under the usual convention that a planar domain is connected and has connectivity \(k\) when its complement in the one-point compactification \(\R^2\cup\{\infty\}\) has exactly \(k\) connected components, set
\[
\mathcal A_k:=\left\{\Omega\subset\R^2:
\begin{array}{l}
\Omega\text{ is nonempty, open, and connected},\\
0<r_\Omega<\infty,\\
(\R^2\cup\{\infty\})\setminus\Omega
\text{ has exactly }k\text{ connected components}
\end{array}\right\}.
\]
This connectivity convention is an explicit provisional interpretation: the supplied excerpt says only ``open multiply connected set of order \(k\)'' and does not define that phrase.

\emph{Source target.} For every \(s\in(0,1)\) and every integer \(k\geq2\), determine the exact value of
\[
T_{s,k}:=\inf_{\Omega\in\mathcal A_k}
\lambda_1^s(\Omega)\,r_\Omega^{2s},
\]
where \(\lambda_1^s(\Omega)\) denotes the fractional first Dirichlet eigenvalue intended by the source. Equivalently, once that eigenvalue and the admissible class are fixed, determine the greatest constant \(T_{s,k}\) such that
\[
\lambda_1^s(\Omega)\,r_\Omega^{2s}\geq T_{s,k}
\qquad\text{for every }\Omega\in\mathcal A_k.
\]
The task does not assert that the infimum is attained.

\emph{Unresolved operator qualification.} The supplied record does not define \(\lambda_1^s\), specify the fractional operator, give its quadratic-form normalization, or state the function space and exterior or boundary conditions. Distinct fractional Dirichlet constructions and normalizations can produce different constants. Consequently, the displayed target remains conditional on recovering those definitions from the source; no particular construction is adopted here.

The original excerpt also refers to an earlier conjecture about the optimal behaviour of \(\vartheta_s\) for \(k\geq2\), but neither \(\vartheta_s\) nor the asserted behaviour is supplied. It further describes the analogous sharp-constant question for the classical Dirichlet Laplacian on simply connected sets, \(k=1\), as a major open problem. These are source-context statements, not additional claims established by this draft.
\subsection{Short English statement}
For each fractional parameter \(0<s<1\) and integer connectivity \(k\geq2\), determine the smallest possible product of the source's first fractional Dirichlet eigenvalue and the \(2s\)-th power of the inradius, over the specified planar open sets with finite inradius. The fractional operator, its normalization, and the precise connectivity convention require confirmation. The supplied record labels this problem open; this statement-only draft requires human mathematical review.
\subsection{Sources}
\begin{itemize}
\item Supplied local UnsolvedMath registry record, permanent ID \(30005294\), title \emph{Sharp Fractional Inradius Constants with Fixed Connectivity}, source problem code \texttt{catalogue entry}, in the Oberwolfach Reports Open Problems collection. No individual proposer is identified; the supplied proposed year is 2022.
\item Supplied source citation: \emph{Mini-Workshop: A Geometric Fairytale full of Spectral Gaps and Random Fruit} (2023), Oberwolfach Reports, pp.~3051--3092, DOI \texttt{10.4171/owr/2022/53}. \url{https://doi.org/10.4171/owr/2022/53}. The supplied extracted passage carries page number 3061. The citation's year 2023 and the record's proposed year 2022 are retained without reconciliation.
\item Supplied literature reference: Bianchi--Brasco, \emph{Optimal lower bound with topological constraints}. \url{https://arxiv.org/abs/2301.08017}.
\item Supplied literature reference: Bianchi--Brasco, \emph{Fractional Makai--Hayman inequality}. \url{https://arxiv.org/abs/2110.11748}.
\item Supplied literature reference: EMS Press, OWR 2022/53. \url{https://doi.org/10.4171/owr/2022/53}.
\item Literature assessment dated exactly \texttt{2026-08-22}: the supplied curation labels the problem open and reports that Bianchi and Brasco obtain explicit lower bounds with sharp parameter ranges and important asymptotics, while the sharp constant for the multiply connected fractional shape-optimization problem remains open. This is supplied dated curation, not an independently verified literature assessment.
\item The registry labels the statement \texttt{corrected\_verified} and reports a comparison with the official Oberwolfach report. That verification is supplied metadata; neither the report nor the linked papers has been independently consulted for this draft.
\end{itemize}
\end{document}
